\chapter{Conclusions and Outlook}
\label{ch:conc}

\section{Results}
The project set out to make the CAN interface between a Stellar-E inverter
and a dynamometer bench automatable. That sentence unpacks into three
artefacts that now agree with each other.

\paragraph{The SR5E1 CAN library is a contract, not a collection of
handlers.}
\texttt{EXECUTE\_COMMAND} commits to a single Motor Control command per
frame. Independent $I_d$ and $I_q$ references exist as first-class codes.
The torque ramp is reachable from CAN, and a flux ramp that the SDK did
not previously offer is implemented through the UI, MCI and STC layers
with the same Q16 increment law as the torque path. PI gains travel as a
packed 32-bit word, numerator and $2^n$ divisor together. ACK/NACK and
\texttt{GET\_REG} no longer share an identifier. Each of those statements
was confirmed on hardware with PCAN-View and TRACE32, not only in
review.

\paragraph{The DBC is a projection of that contract.}
Identifiers follow the $+$\texttt{0x30}/$+$\texttt{0x40} offset rule,
with an explicit exception for the independent current commands where
the naive offset would have collided. Acknowledgement codes
\texttt{0xF0} and \texttt{0xFF} are payload values of message
\texttt{0x10}, not identifiers. Signal layouts---bit-masks, little-endian
physical values, packed PI words---match the C unpacking. The
SET\_REG discrepancy against the dishwasher specification was resolved
by documenting the firmware as authoritative and encoding that layout
in the DBC, rather than by breaking boards already in the field.

\paragraph{An operator can drive the inverter without composing
hexadecimal.}
Vector CANalyzer owns timing and encoding. A Python GUI, talking only
to system variables through the COM API, exposes start/stop, speed
ramps, three PI loops, independent current axes and live plots. CAPL
ready-flags keep incomplete multi-parameter commands off the bus.

Taken together, the chain supports the use-cases that motivated the
work: a speed sweep with a defined duration, a current-loop tuning
session with live $I_d$/$I_q$ charts, a PI write that round-trips, and
a diagnostic stream at \SI{100}{\milli\second} that a CAPL script or
the dyno computer can consume without a human in the loop.

\section{Limitations}
Several limitations are accepted rather than accidental.

\begin{itemize}[leftmargin=1.6em]
  \item The physical layer was validated as CAN~2.0 with 11-bit
        identifiers. The FDCAN peripheral can do more; the DBC and the
        GUI cannot, until someone extends both.
  \item Range checking of current and speed set-points is left to the
        Motor Control SDK. A hostile or mistaken CAPL script can still
        request a value that the regulators will saturate. A future
        saturation layer on the DYNO2DW boundary would make NACKs more
        informative, at the cost of a second source of limits.
  \item The GUI is a Windows process that requires a licensed CANalyzer
        installation. It is the right tool for this laboratory and the
        wrong tool for a headless continuous-integration rack. The
        sysvar/DBC contract, however, is reusable from a CAPL-only
        configuration or from any other COM client.
  \item Functional-safety aspects of Stellar-E (lockstep, watchdog,
        ISO~26262 assumptions) were out of scope. The interface is a
        laboratory automation layer, not a safety element.
\end{itemize}

\section{Outlook}
Three extensions follow naturally from the present contract.

\paragraph{Scripted characterisation.}
With the DBC stable, a CAPL (or Python-COM) script can replay a speed
and current profile, log \texttt{REG\_VALUE} and the diagnostic stack,
and produce a characterisation table without the GUI. The ready-flag
discipline is exactly the guard that such a script needs.

\paragraph{CAN FD and higher-rate telemetry.}
Ten 4-byte messages every \SI{100}{\milli\second} are enough for
operator gauges and tight for a fast current loop. Packing several
registers into a single CAN-FD 64-byte frame, or shortening the period
for a subset of signals, is a DBC-compatible change once the TX path
is measured not to drop stacks under \texttt{FCP\_TRANSFER\_IDLE}
back-pressure.

\paragraph{Bidirectional PI scheduling.}
The packed write of \cref{sec:pi-pack} is atomic per gain. A single
message that writes $K_p$ and $K_i$ of one loop together, or a
read-modify-write of a loop's complete state, would make automated
tuning (relay, Ziegler--Nichols, iterative feedback tuning) less
sensitive to a control period that lands between two CAN frames.

\section{Closing remark}
The technical content of this thesis is a handful of C cases, a DBC
file and a Python window. The engineering content is the refusal to let
those three disagree. A bench that is automated on top of a silently
wrong identifier map is worse than a bench that is operated by hand:
it produces plots that look like measurements and are not. The
debug chain of \cref{ch:val}---inject, halt, step, compare---is the
method that keeps the plots honest, and it is the method this
laboratory can reuse the next time a protocol code is added.
